% GENERATED FILE. Edit canonical lessons, then run npm run generate:books.
% canonical-source-hash: fnv1a64:e2d88f537513e035
% canonical-lessons: PT-C457-carvalho, PT-C457-pinheiro, PT-C457-sobreiro, PT-C457-cortica, PT-C457-podar

\chapter{As árvores}
\label{ch:pt-as-arvores}
\hlchaptermodality{457}

\begin{chapteropening}
\textbf{By the end of this chapter:} \emph{I can name trees and care for them: the oak, the pine, the cork oak, cork and pruning.}

Complete the last lesson of chapter 457: I can name trees and care for them: the oak, the pine, the cork oak, cork and pruning.
\end{chapteropening}

\section[o carvalho]{o carvalho — the oak (tree)}

\label{lesson:PT-C457-carvalho}



\begin{quote}
Before the new one: say the Portuguese for the bush, then the Portuguese for to wade across.
\end{quote}

\subsection*{You'll want to know: o carvalho}
\textbf{o carvalho} — \textquotedblleft{}the oak (tree)\textquotedblright{}.

\begin{quote}
No parque há um carvalho com duzentos anos.
\end{quote}

\begin{quote}
\emph{In the park there is a two-hundred-year-old oak.}
\end{quote}

\emph{Madeira de carvalho} — oak wood.

\begin{cousinweb}
\emph{carvalho} probably comes from a pre-Roman language of the Iberian Peninsula, but the origin is not settled; EtymDB lists it without a source.
\end{cousinweb}

\begin{grammarlens}[title={What you've built}]
One more word for telling, explaining and describing at B1.
\end{grammarlens}

\subsection*{Your turn}
\begin{itemize}
\raggedright
  \item \emph{Say it:} \emph{o carvalho}
  \item \emph{Say it:} \emph{o carvalho}, once more
  \item \emph{Say it:} a sentence of your own with \emph{o carvalho}
  \item \emph{Recall:} say the Portuguese for the bush, then the Portuguese for to wade across, then say \emph{o carvalho} again
\end{itemize}

\begin{tcolorbox}[breakable,colback=teal!4,colframe=teal!35!black,title={Before you move on}]
What does o carvalho mean? (\textquotedblleft{}The oak (tree)\textquotedblright{}.) Say it once more.
\end{tcolorbox}
\section[o pinheiro]{o pinheiro — the pine (tree)}

\label{lesson:PT-C457-pinheiro}



\begin{quote}
Before the new one: say the Portuguese for to wade across, then the Portuguese for the oak.
\end{quote}

\subsection*{You'll want to know: o pinheiro}
\textbf{o pinheiro} — \textquotedblleft{}the pine (tree)\textquotedblright{}.

\begin{quote}
Os pinheiros cheiram bem depois da chuva.
\end{quote}

\begin{quote}
\emph{The pine trees smell good after the rain.}
\end{quote}

\emph{A pinha} — the pine cone; \emph{o pinhão} — the pine nut.

\begin{cousinweb}
\emph{pinheiro} is built on \emph{pinho} 'pine wood' (EtymDB) with \emph{-eiro}; \emph{pinho} is inherited from Latin \emph{pīnus} 'pine'.
\end{cousinweb}

\begin{grammarlens}[title={What you've built}]
One more word for telling, explaining and describing at B1.
\end{grammarlens}

\subsection*{Your turn}
\begin{itemize}
\raggedright
  \item \emph{Say it:} \emph{o pinheiro}
  \item \emph{Say it:} \emph{o pinheiro}, once more
  \item \emph{Say it:} a sentence of your own with \emph{o pinheiro}
  \item \emph{Recall:} say the Portuguese for to wade across, then the Portuguese for the oak, then say \emph{o pinheiro} again
\end{itemize}

\begin{tcolorbox}[breakable,colback=teal!4,colframe=teal!35!black,title={Before you move on}]
What does o pinheiro mean? (\textquotedblleft{}The pine (tree)\textquotedblright{}.) Say it once more.
\end{tcolorbox}
\section[o sobreiro]{o sobreiro — the cork oak}

\label{lesson:PT-C457-sobreiro}



\begin{quote}
Before the new one: say the Portuguese for the oak, then the Portuguese for the pine.
\end{quote}

\subsection*{You'll want to know: o sobreiro}
\textbf{o sobreiro} — \textquotedblleft{}the cork oak\textquotedblright{}.

\begin{quote}
O sobreiro cresce bem no sul do país.
\end{quote}

\begin{quote}
\emph{The cork oak grows well in the south of the country.}
\end{quote}

It is a protected tree in Portugal.

\begin{cousinweb}
\emph{sobreiro} is built on old \emph{sobro} (EtymDB) with \emph{-eiro}; \emph{sobro} is inherited from Latin \emph{sūber} 'cork oak' (Lewis \& Short).
\end{cousinweb}

\begin{grammarlens}[title={What you've built}]
One more word for telling, explaining and describing at B1.
\end{grammarlens}

\subsection*{Your turn}
\begin{itemize}
\raggedright
  \item \emph{Say it:} \emph{o sobreiro}
  \item \emph{Say it:} \emph{o sobreiro}, once more
  \item \emph{Say it:} a sentence of your own with \emph{o sobreiro}
  \item \emph{Recall:} say the Portuguese for the oak, then the Portuguese for the pine, then say \emph{o sobreiro} again
\end{itemize}

\begin{tcolorbox}[breakable,colback=teal!4,colframe=teal!35!black,title={Before you move on}]
What does o sobreiro mean? (\textquotedblleft{}The cork oak\textquotedblright{}.) Say it once more.
\end{tcolorbox}
\section[a cortiça]{a cortiça — cork}

\label{lesson:PT-C457-cortica}



\begin{quote}
Before the new one: say the Portuguese for the pine, then the Portuguese for the cork oak.
\end{quote}

\subsection*{You'll want to know: a cortiça}
\textbf{a cortiça} — \textquotedblleft{}cork\textquotedblright{}.

\begin{quote}
As rolhas das garrafas são de cortiça.
\end{quote}

\begin{quote}
\emph{Bottle corks are made of cork.}
\end{quote}

Bark cut from the cork oak every nine years.

\begin{cousinweb}
\emph{cortiça} is inherited from Latin \emph{corticea} 'made of bark' (EtymDB), from \emph{cortex} 'bark'.
\end{cousinweb}

\begin{grammarlens}[title={What you've built}]
One more word for telling, explaining and describing at B1.
\end{grammarlens}

\subsection*{Your turn}
\begin{itemize}
\raggedright
  \item \emph{Say it:} \emph{a cortiça}
  \item \emph{Say it:} \emph{a cortiça}, once more
  \item \emph{Say it:} a sentence of your own with \emph{a cortiça}
  \item \emph{Recall:} say the Portuguese for the pine, then the Portuguese for the cork oak, then say \emph{a cortiça} again
\end{itemize}

\begin{tcolorbox}[breakable,colback=teal!4,colframe=teal!35!black,title={Before you move on}]
What does a cortiça mean? (\textquotedblleft{}Cork\textquotedblright{}.) Say it once more.
\end{tcolorbox}
\section[podar]{podar — to prune}

\label{lesson:PT-C457-podar}



\begin{quote}
Before the new one: say the Portuguese for the cork oak, then the Portuguese for cork.
\end{quote}

\subsection*{You'll want to know: podar}
\textbf{podar} — \textquotedblleft{}to prune\textquotedblright{}.

\begin{quote}
O jardineiro poda as árvores no inverno.
\end{quote}

\begin{quote}
\emph{The gardener prunes the trees in winter.}
\end{quote}

\emph{A poda} — the pruning.

\begin{cousinweb}
\emph{podar} is inherited from Latin \emph{putāre} 'to cut back branches; to reckon' (EtymDB).
\end{cousinweb}

\begin{grammarlens}[title={What you've built}]
One more word for telling, explaining and describing at B1.
\end{grammarlens}

\subsection*{Your turn}
\begin{itemize}
\raggedright
  \item \emph{Say it:} \emph{podar}
  \item \emph{Say it:} \emph{podar}, once more
  \item \emph{Say it:} a sentence of your own with \emph{podar}
  \item \emph{Recall:} say the Portuguese for the cork oak, then the Portuguese for cork, then say \emph{podar} again
\end{itemize}

\begin{tcolorbox}[breakable,colback=teal!4,colframe=teal!35!black,title={Before you move on}]
What does podar mean? (\textquotedblleft{}To prune\textquotedblright{}.) Say it once more.
\end{tcolorbox}
